\documentclass[aspectratio=169,10pt]{beamer}
\usetheme{Madrid}
\usecolortheme{default}
\usepackage[utf8]{inputenc}
\usepackage[T1]{fontenc}
\usepackage{lmodern}
\usepackage{amsmath,amssymb}
\usepackage{enumitem}
\usepackage{xcolor}

\setbeamertemplate{navigation symbols}{}
\setbeamertemplate{footline}[frame number]
\setbeamertemplate{itemize items}[circle]
\setbeamerfont{frametitle}{size=\large,series=\bfseries}

\title{CIA Part 1 -- Governance, Risk Management, and Control}
\subtitle{100 Hard Questions with Detailed Explanations}
\author{Study Deck}
\date{\today}

\begin{document}
	
	\begin{frame}
		\titlepage
	\end{frame}
	
	\begin{frame}{Format \& Coverage}
		\small
		\begin{itemize}
			\item 100 hard questions on Governance, Risk Management, and Control.
			\item Each question is followed by a detailed explanation slide.
			\item Coverage: governance structures, board/audit committee roles,
			ERM frameworks (COSO ERM, ISO 31000), risk appetite/tolerance,
			inherent/residual/control/detection risk, COSO Internal Control,
			control types, control deficiencies, material weaknesses,
			fraud risk, IT general controls, and governance failures.
		\end{itemize}
	\end{frame}
	
	% =========================================================
	% Q1
	% =========================================================
	\begin{frame}[t]{Q1}
		\small
		\textbf{Question:} Which of the following best describes the primary responsibility for governance in an organization?
		\begin{enumerate}[label=\Alph*.,itemsep=1pt,topsep=2pt]
			\item Management
			\item The board of directors
			\item Internal audit
			\item External audit
		\end{enumerate}
		\vspace{4pt}
		\textbf{Correct Answer: B}
	\end{frame}
	
	\begin{frame}[t]{Q1 -- Explanation}
		\small
		\textbf{Correct: B.} The board holds ultimate responsibility for governance, setting direction, oversight, and accountability. \\[4pt]
		\textbf{A wrong:} Management executes governance but does not own it. \\[2pt]
		\textbf{C wrong:} Internal audit provides assurance on governance. \\[2pt]
		\textbf{D wrong:} External audit provides independent opinion on financial statements.
	\end{frame}
	
	% =========================================================
	% Q2
	% =========================================================
	\begin{frame}[t]{Q2}
		\small
		\textbf{Question:} Which governance principle is most directly supported by having a majority of independent directors?
		\begin{enumerate}[label=\Alph*.,itemsep=1pt,topsep=2pt]
			\item Transparency
			\item Accountability
			\item Objectivity in oversight
			\item Profit maximization
		\end{enumerate}
		\vspace{4pt}
		\textbf{Correct Answer: C}
	\end{frame}
	
	\begin{frame}[t]{Q2 -- Explanation}
		\small
		\textbf{Correct: C.} Independent directors reduce conflicts of interest and support objective oversight. \\[4pt]
		\textbf{A wrong:} Transparency is about disclosure. \\[2pt]
		\textbf{B wrong:} Accountability is about answerability. \\[2pt]
		\textbf{D wrong:} Profit is not a governance principle.
	\end{frame}
	
	% =========================================================
	% Q3
	% =========================================================
	\begin{frame}[t]{Q3}
		\small
		\textbf{Question:} The audit committee's primary role in governance is to:
		\begin{enumerate}[label=\Alph*.,itemsep=1pt,topsep=2pt]
			\item Manage daily operations.
			\item Oversee financial reporting, internal control, and audit functions.
			\item Set product strategy.
			\item Hire all staff.
		\end{enumerate}
		\vspace{4pt}
		\textbf{Correct Answer: B}
	\end{frame}
	
	\begin{frame}[t]{Q3 -- Explanation}
		\small
		\textbf{Correct: B.} The audit committee oversees financial reporting, internal control, internal audit, and external audit. \\[4pt]
		\textbf{A wrong:} Operations are management's job. \\[2pt]
		\textbf{C wrong:} Strategy is the full board/management. \\[2pt]
		\textbf{D wrong:} Hiring is management/HR.
	\end{frame}
	
	% =========================================================
	% Q4
	% =========================================================
	\begin{frame}[t]{Q4}
		\small
		\textbf{Question:} Which of the following is the best evidence of effective governance?
		\begin{enumerate}[label=\Alph*.,itemsep=1pt,topsep=2pt]
			\item High profit margins.
			\item Independent board oversight and effective risk management.
			\item Large market share.
			\item Low employee turnover.
		\end{enumerate}
		\vspace{4pt}
		\textbf{Correct Answer: B}
	\end{frame}
	
	\begin{frame}[t]{Q4 -- Explanation}
		\small
		\textbf{Correct: B.} Effective governance is demonstrated by independent oversight, accountability, and risk management. \\[4pt]
		\textbf{A wrong:} Profit is a result, not governance evidence. \\[2pt]
		\textbf{C wrong:} Market share is a performance metric. \\[2pt]
		\textbf{D wrong:} Turnover is HR-related.
	\end{frame}
	
	% =========================================================
	% Q5
	% =========================================================
	\begin{frame}[t]{Q5}
		\small
		\textbf{Question:} Which is NOT a principle of the OECD governance framework?
		\begin{enumerate}[label=\Alph*.,itemsep=1pt,topsep=2pt]
			\item Ensuring the basis for an effective governance framework.
			\item The rights of shareholders.
			\item Equitable treatment of shareholders.
			\item Maximizing executive bonuses.
		\end{enumerate}
		\vspace{4pt}
		\textbf{Correct Answer: D}
	\end{frame}
	
	\begin{frame}[t]{Q5 -- Explanation}
		\small
		\textbf{Correct: D.} OECD principles cover shareholder rights, equitable treatment, disclosure, board responsibilities, etc., not executive bonuses. \\[4pt]
		\textbf{A wrong:} That is an OECD principle. \\[2pt]
		\textbf{B wrong:} That is an OECD principle. \\[2pt]
		\textbf{C wrong:} That is an OECD principle.
	\end{frame}
	
	% =========================================================
	% Q6
	% =========================================================
	\begin{frame}[t]{Q6}
		\small
		\textbf{Question:} Which best describes the difference between governance and management?
		\begin{enumerate}[label=\Alph*.,itemsep=1pt,topsep=2pt]
			\item Governance sets direction and oversight; management executes.
			\item Governance executes; management sets direction.
			\item They are the same.
			\item Governance is only for public companies.
		\end{enumerate}
		\vspace{4pt}
		\textbf{Correct Answer: A}
	\end{frame}
	
	\begin{frame}[t]{Q6 -- Explanation}
		\small
		\textbf{Correct: A.} Governance sets direction, oversight, and accountability; management runs day-to-day operations. \\[4pt]
		\textbf{B wrong:} Roles are reversed. \\[2pt]
		\textbf{C wrong:} They are distinct. \\[2pt]
		\textbf{D wrong:} All organizations need governance.
	\end{frame}
	
	% =========================================================
	% Q7
	% =========================================================
	\begin{frame}[t]{Q7}
		\small
		\textbf{Question:} Which of the following is the best description of "enterprise risk management" (ERM)?
		\begin{enumerate}[label=\Alph*.,itemsep=1pt,topsep=2pt]
			\item A process to identify, assess, and manage risks across the organization.
			\item A one-time risk assessment.
			\item A compliance checklist.
			\item A financial audit.
		\end{enumerate}
		\vspace{4pt}
		\textbf{Correct Answer: A}
	\end{frame}
	
	\begin{frame}[t]{Q7 -- Explanation}
		\small
		\textbf{Correct: A.} ERM is a structured, ongoing process to identify, assess, respond to, and monitor risks across the enterprise. \\[4pt]
		\textbf{B wrong:} ERM is ongoing. \\[2pt]
		\textbf{C wrong:} It is broader than compliance. \\[2pt]
		\textbf{D wrong:} Financial audit is assurance.
	\end{frame}
	
	% =========================================================
	% Q8
	% =========================================================
	\begin{frame}[t]{Q8}
		\small
		\textbf{Question:} In COSO ERM 2017, which is NOT one of the five components?
		\begin{enumerate}[label=\Alph*.,itemsep=1pt,topsep=2pt]
			\item Governance and culture.
			\item Strategy and objective-setting.
			\item Performance.
			\item Financial reporting.
		\end{enumerate}
		\vspace{4pt}
		\textbf{Correct Answer: D}
	\end{frame}
	
	\begin{frame}[t]{Q8 -- Explanation}
		\small
		\textbf{Correct: D.} COSO ERM 2017 components: Governance \& Culture, Strategy \& Objective-Setting, Performance, Review \& Revision, Information, Communication \& Reporting. \\[4pt]
		\textbf{A wrong:} That is a component. \\[2pt]
		\textbf{B wrong:} That is a component. \\[2pt]
		\textbf{C wrong:} That is a component.
	\end{frame}
	
	% =========================================================
	% Q9
	% =========================================================
	\begin{frame}[t]{Q9}
		\small
		\textbf{Question:} Which best describes the difference between risk appetite and risk tolerance?
		\begin{enumerate}[label=\Alph*.,itemsep=1pt,topsep=2pt]
			\item Appetite is the amount of risk willing to accept; tolerance is the acceptable variation around it.
			\item They are the same.
			\item Tolerance is set by the auditor.
			\item Appetite is set by the auditor.
		\end{enumerate}
		\vspace{4pt}
		\textbf{Correct Answer: A}
	\end{frame}
	
	\begin{frame}[t]{Q9 -- Explanation}
		\small
		\textbf{Correct: A.} Risk appetite is the amount/type of risk an entity is willing to accept; risk tolerance is the acceptable variation around that appetite. \\[4pt]
		\textbf{B wrong:} They are distinct. \\[2pt]
		\textbf{C wrong:} Tolerance is set by management/board. \\[2pt]
		\textbf{D wrong:} Appetite is set by the board.
	\end{frame}
	
	% =========================================================
	% Q10
	% =========================================================
	\begin{frame}[t]{Q10}
		\small
		\textbf{Question:} Which of the following best describes "residual risk"?
		\begin{enumerate}[label=\Alph*.,itemsep=1pt,topsep=2pt]
			\item Risk before controls.
			\item Risk remaining after controls.
			\item Risk of fraud.
			\item Risk of auditor error.
		\end{enumerate}
		\vspace{4pt}
		\textbf{Correct Answer: B}
	\end{frame}
	
	\begin{frame}[t]{Q10 -- Explanation}
		\small
		\textbf{Correct: B.} Residual risk = inherent risk minus the effect of controls. \\[4pt]
		\textbf{A wrong:} That is inherent risk. \\[2pt]
		\textbf{C wrong:} Fraud is one type of risk. \\[2pt]
		\textbf{D wrong:} Auditor error is detection risk.
	\end{frame}
	
	% =========================================================
	% Q11
	% =========================================================
	\begin{frame}[t]{Q11}
		\small
		\textbf{Question:} Which is the best description of "control risk"?
		\begin{enumerate}[label=\Alph*.,itemsep=1pt,topsep=2pt]
			\item Risk that a control fails to prevent or detect a misstatement.
			\item Risk before controls.
			\item Risk of fraud.
			\item Risk of market loss.
		\end{enumerate}
		\vspace{4pt}
		\textbf{Correct Answer: A}
	\end{frame}
	
	\begin{frame}[t]{Q11 -- Explanation}
		\small
		\textbf{Correct: A.} Control risk is the risk that a control will not prevent or detect a material misstatement on a timely basis. \\[4pt]
		\textbf{B wrong:} That is inherent risk. \\[2pt]
		\textbf{C wrong:} Fraud is a source. \\[2pt]
		\textbf{D wrong:} Market loss is market risk.
	\end{frame}
	
	% =========================================================
	% Q12
	% =========================================================
	\begin{frame}[t]{Q12}
		\small
		\textbf{Question:} Which is the best description of "detection risk"?
		\begin{enumerate}[label=\Alph*.,itemsep=1pt,topsep=2pt]
			\item Risk that audit procedures fail to detect a misstatement.
			\item Risk before controls.
			\item Risk of fraud.
			\item Risk of market loss.
		\end{enumerate}
		\vspace{4pt}
		\textbf{Correct Answer: A}
	\end{frame}
	
	\begin{frame}[t]{Q12 -- Explanation}
		\small
		\textbf{Correct: A.} Detection risk is the risk that the auditor's procedures will not detect a material misstatement. \\[4pt]
		\textbf{B wrong:} That is inherent risk. \\[2pt]
		\textbf{C wrong:} Fraud is a source. \\[2pt]
		\textbf{D wrong:} Market loss is market risk.
	\end{frame}
	
	% =========================================================
	% Q13
	% =========================================================
	\begin{frame}[t]{Q13}
		\small
		\textbf{Question:} Which of the following is a preventive control?
		\begin{enumerate}[label=\Alph*.,itemsep=1pt,topsep=2pt]
			\item Bank reconciliation.
			\item Segregation of duties.
			\item Variance analysis.
			\item Exception report.
		\end{enumerate}
		\vspace{4pt}
		\textbf{Correct Answer: B}
	\end{frame}
	
	\begin{frame}[t]{Q13 -- Explanation}
		\small
		\textbf{Correct: B.} Segregation of duties prevents fraud/error. \\[4pt]
		\textbf{A wrong:} Reconciliation is detective. \\[2pt]
		\textbf{C wrong:} Variance analysis is detective. \\[2pt]
		\textbf{D wrong:} Exception reports are detective.
	\end{frame}
	
	% =========================================================
	% Q14
	% =========================================================
	\begin{frame}[t]{Q14}
		\small
		\textbf{Question:} Which of the following is a detective control?
		\begin{enumerate}[label=\Alph*.,itemsep=1pt,topsep=2pt]
			\item Authorization.
			\item Physical safeguards.
			\item Bank reconciliation.
			\item Password policy.
		\end{enumerate}
		\vspace{4pt}
		\textbf{Correct Answer: C}
	\end{frame}
	
	\begin{frame}[t]{Q14 -- Explanation}
		\small
		\textbf{Correct: C.} Reconciliation detects discrepancies after they occur. \\[4pt]
		\textbf{A wrong:} Authorization is preventive. \\[2pt]
		\textbf{B wrong:} Safeguards are preventive. \\[2pt]
		\textbf{D wrong:} Password policy is preventive.
	\end{frame}
	
	% =========================================================
	% Q15
	% =========================================================
	\begin{frame}[t]{Q15}
		\small
		\textbf{Question:} Which of the following is a corrective control?
		\begin{enumerate}[label=\Alph*.,itemsep=1pt,topsep=2pt]
			\item Authorization.
			\item Segregation of duties.
			\item Backup and recovery procedures.
			\item Password policy.
		\end{enumerate}
		\vspace{4pt}
		\textbf{Correct Answer: C}
	\end{frame}
	
	\begin{frame}[t]{Q15 -- Explanation}
		\small
		\textbf{Correct: C.} Corrective controls restore systems/processes after an incident (e.g., backup recovery). \\[4pt]
		\textbf{A wrong:} Authorization is preventive. \\[2pt]
		\textbf{B wrong:} Segregation is preventive. \\[2pt]
		\textbf{D wrong:} Password policy is preventive.
	\end{frame}
	
	% =========================================================
	% Q16
	% =========================================================
	\begin{frame}[t]{Q16}
		\small
		\textbf{Question:} Which is the best description of "inherent risk"?
		\begin{enumerate}[label=\Alph*.,itemsep=1pt,topsep=2pt]
			\item Risk before considering controls.
			\item Risk after controls.
			\item Risk of fraud.
			\item Risk of market loss.
		\end{enumerate}
		\vspace{4pt}
		\textbf{Correct Answer: A}
	\end{frame}
	
	\begin{frame}[t]{Q16 -- Explanation}
		\small
		\textbf{Correct: A.} Inherent risk is the risk of an event occurring before controls are applied. \\[4pt]
		\textbf{B wrong:} That is residual risk. \\[2pt]
		\textbf{C wrong:} Fraud is a source. \\[2pt]
		\textbf{D wrong:} Market loss is market risk.
	\end{frame}
	
	% =========================================================
	% Q17
	% =========================================================
	\begin{frame}[t]{Q17}
		\small
		\textbf{Question:} Which of the following is NOT a component of the COSO Internal Control -- Integrated Framework?
		\begin{enumerate}[label=\Alph*.,itemsep=1pt,topsep=2pt]
			\item Control environment.
			\item Risk assessment.
			\item Marketing strategy.
			\item Monitoring.
		\end{enumerate}
		\vspace{4pt}
		\textbf{Correct Answer: C}
	\end{frame}
	
	\begin{frame}[t]{Q17 -- Explanation}
		\small
		\textbf{Correct: C.} COSO's five components: Control Environment, Risk Assessment, Control Activities, Information \& Communication, Monitoring. \\[4pt]
		\textbf{A wrong:} Control environment is COSO. \\[2pt]
		\textbf{B wrong:} Risk assessment is COSO. \\[2pt]
		\textbf{D wrong:} Monitoring is COSO.
	\end{frame}
	
	% =========================================================
	% Q18
	% =========================================================
	\begin{frame}[t]{Q18}
		\small
		\textbf{Question:} Which best describes "tone at the top"?
		\begin{enumerate}[label=\Alph*.,itemsep=1pt,topsep=2pt]
			\item The CEO's salary.
			\item Management's commitment to ethics and control.
			\item The marketing slogan.
			\item The IT password policy.
		\end{enumerate}
		\vspace{4pt}
		\textbf{Correct Answer: B}
	\end{frame}
	
	\begin{frame}[t]{Q18 -- Explanation}
		\small
		\textbf{Correct: B.} Tone at the top is the control environment element reflecting leadership's ethics and commitment to control. \\[4pt]
		\textbf{A wrong:} Compensation is unrelated. \\[2pt]
		\textbf{C wrong:} Marketing slogans are not control environment. \\[2pt]
		\textbf{D wrong:} Password policy is one IT control.
	\end{frame}
	
	% =========================================================
	% Q19
	% =========================================================
	\begin{frame}[t]{Q19}
		\small
		\textbf{Question:} Which is a limitation of internal control?
		\begin{enumerate}[label=\Alph*.,itemsep=1pt,topsep=2pt]
			\item Cost-benefit considerations.
			\item Segregation of duties.
			\item Physical safeguards.
			\item Authorization procedures.
		\end{enumerate}
		\vspace{4pt}
		\textbf{Correct Answer: A}
	\end{frame}
	
	\begin{frame}[t]{Q19 -- Explanation}
		\small
		\textbf{Correct: A.} Controls are limited by cost-benefit, human error, collusion, and management override. \\[4pt]
		\textbf{B wrong:} Segregation is a control. \\[2pt]
		\textbf{C wrong:} Safeguards are controls. \\[2pt]
		\textbf{D wrong:} Authorization is a control.
	\end{frame}
	
	% =========================================================
	% Q20
	% =========================================================
	\begin{frame}[t]{Q20}
		\small
		\textbf{Question:} Which is the best description of a "control deficiency"?
		\begin{enumerate}[label=\Alph*.,itemsep=1pt,topsep=2pt]
			\item Design or operation does not prevent/detect misstatements timely.
			\item A documented control.
			\item A management disagreement.
			\item Zero risk.
		\end{enumerate}
		\vspace{4pt}
		\textbf{Correct Answer: A}
	\end{frame}
	
	\begin{frame}[t]{Q20 -- Explanation}
		\small
		\textbf{Correct: A.} A deficiency exists when a control is missing, poorly designed, or not operating so misstatements are not prevented/detected timely. \\[4pt]
		\textbf{B wrong:} Documentation is a strength. \\[2pt]
		\textbf{C wrong:} Disagreement doesn't create a deficiency. \\[2pt]
		\textbf{D wrong:} Zero risk is unrealistic.
	\end{frame}
	
	% =========================================================
	% Q21
	% =========================================================
	\begin{frame}[t]{Q21}
		\small
		\textbf{Question:} Which is the best description of a "material weakness"?
		\begin{enumerate}[label=\Alph*.,itemsep=1pt,topsep=2pt]
			\item A minor deficiency.
			\item A deficiency creating a reasonable possibility of material misstatement.
			\item A documentation gap.
			\item A training issue.
		\end{enumerate}
		\vspace{4pt}
		\textbf{Correct Answer: B}
	\end{frame}
	
	\begin{frame}[t]{Q21 -- Explanation}
		\small
		\textbf{Correct: B.} A material weakness creates a reasonable possibility that a material misstatement will not be prevented or detected timely. \\[4pt]
		\textbf{A wrong:} Minor deficiencies are less severe. \\[2pt]
		\textbf{C wrong:} Documentation gaps may be deficiencies but not automatically material. \\[2pt]
		\textbf{D wrong:} Training issues may contribute but are not the definition.
	\end{frame}
	
	% =========================================================
	% Q22
	% =========================================================
	\begin{frame}[t]{Q22}
		\small
		\textbf{Question:} Which is the best description of a "significant deficiency"?
		\begin{enumerate}[label=\Alph*.,itemsep=1pt,topsep=2pt]
			\item Less severe than a material weakness but important enough to report.
			\item Always material.
			\item A minor clerical error.
			\item Only relevant to IT.
		\end{enumerate}
		\vspace{4pt}
		\textbf{Correct Answer: A}
	\end{frame}
	
	\begin{frame}[t]{Q22 -- Explanation}
		\small
		\textbf{Correct: A.} A significant deficiency is less severe than a material weakness but important enough to bring to the attention of those charged with governance. \\[4pt]
		\textbf{B wrong:} It is less severe than material. \\[2pt]
		\textbf{C wrong:} Minor clerical errors are not significant deficiencies. \\[2pt]
		\textbf{D wrong:} It applies beyond IT.
	\end{frame}
	
	% =========================================================
	% Q23
	% =========================================================
	\begin{frame}[t]{Q23}
		\small
		\textbf{Question:} Which is the best description of "fraud risk"?
		\begin{enumerate}[label=\Alph*.,itemsep=1pt,topsep=2pt]
			\item Risk of intentional misstatement or theft.
			\item Risk of human error.
			\item Risk of market loss.
			\item Risk of IT failure.
		\end{enumerate}
		\vspace{4pt}
		\textbf{Correct Answer: A}
	\end{frame}
	
	\begin{frame}[t]{Q23 -- Explanation}
		\small
		\textbf{Correct: A.} Fraud risk is the risk of intentional acts (misstatement, theft) by individuals. \\[4pt]
		\textbf{B wrong:} Human error is not fraud. \\[2pt]
		\textbf{C wrong:} Market loss is market risk. \\[2pt]
		\textbf{D wrong:} IT failure is operational risk.
	\end{frame}
	
	% =========================================================
	% Q24
	% =========================================================
	\begin{frame}[t]{Q24}
		\small
		\textbf{Question:} The fraud triangle consists of:
		\begin{enumerate}[label=\Alph*.,itemsep=1pt,topsep=2pt]
			\item Pressure, opportunity, rationalization.
			\item Risk, control, detection.
			\item Plan, do, check.
			\item Cash, inventory, receivables.
		\end{enumerate}
		\vspace{4pt}
		\textbf{Correct Answer: A}
	\end{frame}
	
	\begin{frame}[t]{Q24 -- Explanation}
		\small
		\textbf{Correct: A.} Cressey's fraud triangle: pressure (incentive), opportunity, and rationalization. \\[4pt]
		\textbf{B wrong:} That's the audit risk model. \\[2pt]
		\textbf{C wrong:} That's the Deming/PDCA cycle. \\[2pt]
		\textbf{D wrong:} Those are balance sheet accounts.
	\end{frame}
	
	% =========================================================
	% Q25
	% =========================================================
	\begin{frame}[t]{Q25}
		\small
		\textbf{Question:} Which best describes the IIA's role in ERM?
		\begin{enumerate}[label=\Alph*.,itemsep=1pt,topsep=2pt]
			\item Own the ERM process.
			\item Provide assurance and consulting on ERM.
			\item Set risk appetite.
			\item Approve risk responses.
		\end{enumerate}
		\vspace{4pt}
		\textbf{Correct Answer: B}
	\end{frame}
	
	\begin{frame}[t]{Q25 -- Explanation}
		\small
		\textbf{Correct: B.} Internal audit provides assurance and consulting on ERM but should not own it or set appetite. \\[4pt]
		\textbf{A wrong:} Management owns ERM. \\[2pt]
		\textbf{C wrong:} The board sets risk appetite. \\[2pt]
		\textbf{D wrong:} Management approves responses.
	\end{frame}
	
	% =========================================================
	% Q26
	% =========================================================
	\begin{frame}[t]{Q26}
		\small
		\textbf{Question:} Which is the best description of "risk response"?
		\begin{enumerate}[label=\Alph*.,itemsep=1pt,topsep=2pt]
			\item Avoid, reduce, share, or accept risk.
			\item Ignore risk.
			\item Transfer all risk.
			\item Eliminate all risk.
		\end{enumerate}
		\vspace{4pt}
		\textbf{Correct Answer: A}
	\end{frame}
	
	\begin{frame}[t]{Q26 -- Explanation}
		\small
		\textbf{Correct: A.} The four responses are avoid, reduce (mitigate), share (transfer), and accept. \\[4pt]
		\textbf{B wrong:} Ignoring is not valid. \\[2pt]
		\textbf{C wrong:} Transfer is only one option. \\[2pt]
		\textbf{D wrong:} Elimination is not always possible.
	\end{frame}
	
	% =========================================================
	% Q27
	% =========================================================
	\begin{frame}[t]{Q27}
		\small
		\textbf{Question:} Which best describes "risk register"?
		\begin{enumerate}[label=\Alph*.,itemsep=1pt,topsep=2pt]
			\item A list of identified risks and their responses.
			\item The audit plan.
			\item The marketing plan.
			\item The financial statements.
		\end{enumerate}
		\vspace{4pt}
		\textbf{Correct Answer: A}
	\end{frame}
	
	\begin{frame}[t]{Q27 -- Explanation}
		\small
		\textbf{Correct: A.} A risk register documents identified risks, their assessment, owners, and responses. \\[4pt]
		\textbf{B wrong:} The audit plan is separate. \\[2pt]
		\textbf{C wrong:} The marketing plan is unrelated. \\[2pt]
		\textbf{D wrong:} Financial statements are outputs.
	\end{frame}
	
	% =========================================================
	% Q28
	% =========================================================
	\begin{frame}[t]{Q28}
		\small
		\textbf{Question:} Which best describes "risk assessment" in COSO?
		\begin{enumerate}[label=\Alph*.,itemsep=1pt,topsep=2pt]
			\item Identifying and analyzing risks to objectives.
			\item Preparing financial statements.
			\item Hiring staff.
			\item Auditing the financial statements.
		\end{enumerate}
		\vspace{4pt}
		\textbf{Correct Answer: A}
	\end{frame}
	
	\begin{frame}[t]{Q28 -- Explanation}
		\small
		\textbf{Correct: A.} Risk assessment identifies and analyzes risks to achieving objectives. \\[4pt]
		\textbf{B wrong:} Financial statement prep is management. \\[2pt]
		\textbf{C wrong:} Hiring is a control activity. \\[2pt]
		\textbf{D wrong:} Auditing is assurance.
	\end{frame}
	
	% =========================================================
	% Q29
	% =========================================================
	\begin{frame}[t]{Q29}
		\small
		\textbf{Question:} Which is the best description of "control activities"?
		\begin{enumerate}[label=\Alph*.,itemsep=1pt,topsep=2pt]
			\item Policies and procedures that help ensure management directives are carried out.
			\item The board's meeting schedule.
			\item The audit plan.
			\item The marketing plan.
		\end{enumerate}
		\vspace{4pt}
		\textbf{Correct Answer: A}
	\end{frame}
	
	\begin{frame}[t]{Q29 -- Explanation}
		\small
		\textbf{Correct: A.} Control activities are policies and procedures that help ensure management directives are executed. \\[4pt]
		\textbf{B wrong:} Meeting schedules are not control activities. \\[2pt]
		\textbf{C wrong:} The audit plan is a planning document. \\[2pt]
		\textbf{D wrong:} Marketing plans are unrelated.
	\end{frame}
	
	% =========================================================
	% Q30
	% =========================================================
	\begin{frame}[t]{Q30}
		\small
		\textbf{Question:} Which best describes "monitoring" in COSO?
		\begin{enumerate}[label=\Alph*.,itemsep=1pt,topsep=2pt]
			\item Assessing the quality of internal control over time.
			\item Preparing the budget.
			\item Hiring staff.
			\item Setting the risk appetite.
		\end{enumerate}
		\vspace{4pt}
		\textbf{Correct Answer: A}
	\end{frame}
	
	\begin{frame}[t]{Q30 -- Explanation}
		\small
		\textbf{Correct: A.} Monitoring assesses the quality of internal control performance over time. \\[4pt]
		\textbf{B wrong:} Budgeting is a control activity. \\[2pt]
		\textbf{C wrong:} Hiring is a control activity. \\[2pt]
		\textbf{D wrong:} Risk appetite is set by the board.
	\end{frame}
	
	% =========================================================
	% Q31
	% =========================================================
	\begin{frame}[t]{Q31}
		\small
		\textbf{Question:} Which best describes "information and communication" in COSO?
		\begin{enumerate}[label=\Alph*.,itemsep=1pt,topsep=2pt]
			\item Identifying, capturing, and communicating relevant information.
			\item The marketing campaign.
			\item The IT budget.
			\item The audit report.
		\end{enumerate}
		\vspace{4pt}
		\textbf{Correct Answer: A}
	\end{frame}
	
	\begin{frame}[t]{Q31 -- Explanation}
		\small
		\textbf{Correct: A.} Information and communication ensures relevant information is identified, captured, and communicated timely. \\[4pt]
		\textbf{B wrong:} Marketing is unrelated. \\[2pt]
		\textbf{C wrong:} IT budget is a resource. \\[2pt]
		\textbf{D wrong:} The audit report is an output.
	\end{frame}
	
	% =========================================================
	% Q32
	% =========================================================
	\begin{frame}[t]{Q32}
		\small
		\textbf{Question:} Which best describes "control environment"?
		\begin{enumerate}[label=\Alph*.,itemsep=1pt,topsep=2pt]
			\item The tone of the organization, influencing control consciousness.
			\item The IT infrastructure.
			\item The audit committee.
			\item The external auditor's opinion.
		\end{enumerate}
		\vspace{4pt}
		\textbf{Correct Answer: A}
	\end{frame}
	
	\begin{frame}[t]{Q32 -- Explanation}
		\small
		\textbf{Correct: A.} The control environment is the tone at the top -- integrity, ethical values, and commitment to competence. \\[4pt]
		\textbf{B wrong:} IT infrastructure is part of control activities. \\[2pt]
		\textbf{C wrong:} The audit committee is part of governance. \\[2pt]
		\textbf{D wrong:} External opinion is not the environment.
	\end{frame}
	
	% =========================================================
	% Q33
	% =========================================================
	\begin{frame}[t]{Q33}
		\small
		\textbf{Question:} Which best describes "internal control"?
		\begin{enumerate}[label=\Alph*.,itemsep=1pt,topsep=2pt]
			\item A process designed to provide reasonable assurance regarding objectives.
			\item A guarantee of no fraud.
			\item The external audit opinion.
			\item The marketing plan.
		\end{enumerate}
		\vspace{4pt}
		\textbf{Correct Answer: A}
	\end{frame}
	
	\begin{frame}[t]{Q33 -- Explanation}
		\small
		\textbf{Correct: A.} Internal control is a process designed to provide reasonable assurance regarding operations, reporting, and compliance objectives. \\[4pt]
		\textbf{B wrong:} No control guarantees no fraud. \\[2pt]
		\textbf{C wrong:} External opinion is separate. \\[2pt]
		\textbf{D wrong:} Marketing is unrelated.
	\end{frame}
	
	% =========================================================
	% Q34
	% =========================================================
	\begin{frame}[t]{Q34}
		\small
		\textbf{Question:} Which is the best description of "compliance risk"?
		\begin{enumerate}[label=\Alph*.,itemsep=1pt,topsep=2pt]
			\item Risk of violating laws, regulations, or policies.
			\item Risk of losing market share.
			\item Risk of IT failure.
			\item Risk of fraud.
		\end{enumerate}
		\vspace{4pt}
		\textbf{Correct Answer: A}
	\end{frame}
	
	\begin{frame}[t]{Q34 -- Explanation}
		\small
		\textbf{Correct: A.} Compliance risk is the risk of violating laws, regulations, rules, or policies. \\[4pt]
		\textbf{B wrong:} Market share is strategic risk. \\[2pt]
		\textbf{C wrong:} IT failure is operational/technology risk. \\[2pt]
		\textbf{D wrong:} Fraud is its own category.
	\end{frame}
	
	% =========================================================
	% Q35
	% =========================================================
	\begin{frame}[t]{Q35}
		\small
		\textbf{Question:} Which is the best description of "operational risk"?
		\begin{enumerate}[label=\Alph*.,itemsep=1pt,topsep=2pt]
			\item Risk of loss from inadequate or failed processes, people, or systems.
			\item Risk of losing market share.
			\item Risk of violating laws.
			\item Risk of fraud.
		\end{enumerate}
		\vspace{4pt}
		\textbf{Correct Answer: A}
	\end{frame}
	
	\begin{frame}[t]{Q35 -- Explanation}
		\small
		\textbf{Correct: A.} Operational risk arises from inadequate or failed internal processes, people, systems, or external events. \\[4pt]
		\textbf{B wrong:} Market share is strategic. \\[2pt]
		\textbf{C wrong:} Legal violations are compliance risk. \\[2pt]
		\textbf{D wrong:} Fraud is a subset.
	\end{frame}
	
	% =========================================================
	% Q36
	% =========================================================
	\begin{frame}[t]{Q36}
		\small
		\textbf{Question:} Which is the best description of "strategic risk"?
		\begin{enumerate}[label=\Alph*.,itemsep=1pt,topsep=2pt]
			\item Risk of failing to achieve strategic objectives.
			\item Risk of IT failure.
			\item Risk of fraud.
			\item Risk of violating laws.
		\end{enumerate}
		\vspace{4pt}
		\textbf{Correct Answer: A}
	\end{frame}
	
	\begin{frame}[t]{Q36 -- Explanation}
		\small
		\textbf{Correct: A.} Strategic risk is the risk of failing to achieve strategic objectives. \\[4pt]
		\textbf{B wrong:} IT failure is operational. \\[2pt]
		\textbf{C wrong:} Fraud is its own category. \\[2pt]
		\textbf{D wrong:} Legal violations are compliance.
	\end{frame}
	
	% =========================================================
	% Q37
	% =========================================================
	\begin{frame}[t]{Q37}
		\small
		\textbf{Question:} Which is the best description of "reputational risk"?
		\begin{enumerate}[label=\Alph*.,itemsep=1pt,topsep=2pt]
			\item Risk of damage to the organization's reputation.
			\item Risk of IT failure.
			\item Risk of fraud.
			\item Risk of violating laws.
		\end{enumerate}
		\vspace{4pt}
		\textbf{Correct Answer: A}
	\end{frame}
	
	\begin{frame}[t]{Q37 -- Explanation}
		\small
		\textbf{Correct: A.} Reputational risk is the risk of damage to the organization's image, brand, or public trust. \\[4pt]
		\textbf{B wrong:} IT failure is operational. \\[2pt]
		\textbf{C wrong:} Fraud is a separate category. \\[2pt]
		\textbf{D wrong:} Legal violations are compliance risk.
	\end{frame}
	
	% =========================================================
	% Q38
	% =========================================================
	\begin{frame}[t]{Q38}
		\small
		\textbf{Question:} Which is the best description of "risk treatment"?
		\begin{enumerate}[label=\Alph*.,itemsep=1pt,topsep=2pt]
			\item Selecting and implementing measures to modify risk.
			\item Ignoring risk.
			\item Transferring all risk.
			\item Eliminating all risk.
		\end{enumerate}
		\vspace{4pt}
		\textbf{Correct Answer: A}
	\end{frame}
	
	\begin{frame}[t]{Q38 -- Explanation}
		\small
		\textbf{Correct: A.} Risk treatment involves selecting and implementing measures to modify risk (avoid, reduce, share, accept). \\[4pt]
		\textbf{B wrong:} Ignoring is not treatment. \\[2pt]
		\textbf{C wrong:} Transfer is one option. \\[2pt]
		\textbf{D wrong:} Elimination is not always possible.
	\end{frame}
	
	% =========================================================
	% Q39
	% =========================================================
	\begin{frame}[t]{Q39}
		\small
		\textbf{Question:} Which best describes the "three lines model"?
		\begin{enumerate}[label=\Alph*.,itemsep=1pt,topsep=2pt]
			\item Management, risk/compliance, internal audit.
			\item Board, CEO, CFO.
			\item External audit, internal audit, regulators.
			\item Marketing, sales, operations.
		\end{enumerate}
		\vspace{4pt}
		\textbf{Correct Answer: A}
	\end{frame}
	
	\begin{frame}[t]{Q39 -- Explanation}
		\small
		\textbf{Correct: A.} The Three Lines Model: first line (management/operations), second line (risk/compliance), third line (internal audit). \\[4pt]
		\textbf{B wrong:} Those are roles, not the model. \\[2pt]
		\textbf{C wrong:} External audit is outside the model. \\[2pt]
		\textbf{D wrong:} Those are functions, not the model.
	\end{frame}
	
	% =========================================================
	% Q40
	% =========================================================
	\begin{frame}[t]{Q40}
		\small
		\textbf{Question:} Which is the best description of the first line in the Three Lines Model?
		\begin{enumerate}[label=\Alph*.,itemsep=1pt,topsep=2pt]
			\item Management owns and manages risk.
			\item Risk/compliance oversees risk.
			\item Internal audit provides assurance.
			\item External audit provides opinion.
		\end{enumerate}
		\vspace{4pt}
		\textbf{Correct Answer: A}
	\end{frame}
	
	\begin{frame}[t]{Q40 -- Explanation}
		\small
		\textbf{Correct: A.} The first line is management, which owns and manages risk. \\[4pt]
		\textbf{B wrong:} Risk/compliance is the second line. \\[2pt]
		\textbf{C wrong:} Internal audit is the third line. \\[2pt]
		\textbf{D wrong:} External audit is outside the model.
	\end{frame}
	
	% =========================================================
	% Q41
	% =========================================================
	\begin{frame}[t]{Q41}
		\small
		\textbf{Question:} Which is the best description of the second line in the Three Lines Model?
		\begin{enumerate}[label=\Alph*.,itemsep=1pt,topsep=2pt]
			\item Risk/compliance functions that oversee and challenge risk.
			\item Management owns risk.
			\item Internal audit provides assurance.
			\item External audit provides opinion.
		\end{enumerate}
		\vspace{4pt}
		\textbf{Correct Answer: A}
	\end{frame}
	
	\begin{frame}[t]{Q41 -- Explanation}
		\small
		\textbf{Correct: A.} The second line is risk/compliance, which oversees and challenges risk management. \\[4pt]
		\textbf{B wrong:} Management is the first line. \\[2pt]
		\textbf{C wrong:} Internal audit is the third line. \\[2pt]
		\textbf{D wrong:} External audit is outside the model.
	\end{frame}
	
	% =========================================================
	% Q42
	% =========================================================
	\begin{frame}[t]{Q42}
		\small
		\textbf{Question:} Which is the best description of the third line in the Three Lines Model?
		\begin{enumerate}[label=\Alph*.,itemsep=1pt,topsep=2pt]
			\item Internal audit provides independent assurance.
			\item Management owns risk.
			\item Risk/compliance oversees risk.
			\item External audit provides opinion.
		\end{enumerate}
		\vspace{4pt}
		\textbf{Correct Answer: A}
	\end{frame}
	
	\begin{frame}[t]{Q42 -- Explanation}
		\small
		\textbf{Correct: A.} The third line is internal audit, providing independent assurance. \\[4pt]
		\textbf{B wrong:} Management is the first line. \\[2pt]
		\textbf{C wrong:} Risk/compliance is the second line. \\[2pt]
		\textbf{D wrong:} External audit is outside the model.
	\end{frame}
	
	% =========================================================
	% Q43
	% =========================================================
	\begin{frame}[t]{Q43}
		\small
		\textbf{Question:} Which best describes "governance"?
		\begin{enumerate}[label=\Alph*.,itemsep=1pt,topsep=2pt]
			\item The processes and structures to direct and control the organization.
			\item The audit plan.
			\item The marketing plan.
			\item The financial statements.
		\end{enumerate}
		\vspace{4pt}
		\textbf{Correct Answer: A}
	\end{frame}
	
	\begin{frame}[t]{Q43 -- Explanation}
		\small
		\textbf{Correct: A.} Governance encompasses the processes and structures used to direct, control, and hold the organization accountable. \\[4pt]
		\textbf{B wrong:} The audit plan is a tool. \\[2pt]
		\textbf{C wrong:} Marketing is unrelated. \\[2pt]
		\textbf{D wrong:} Financial statements are outputs.
	\end{frame}
	
	% =========================================================
	% Q44
	% =========================================================
	\begin{frame}[t]{Q44}
		\small
		\textbf{Question:} Which is the best description of "board oversight"?
		\begin{enumerate}[label=\Alph*.,itemsep=1pt,topsep=2pt]
			\item Monitoring management and ensuring accountability.
			\item Managing daily operations.
			\item Preparing workpapers.
			\item Hiring all staff.
		\end{enumerate}
		\vspace{4pt}
		\textbf{Correct Answer: A}
	\end{frame}
	
	\begin{frame}[t]{Q44 -- Explanation}
		\small
		\textbf{Correct: A.} Board oversight monitors management and ensures accountability to stakeholders. \\[4pt]
		\textbf{B wrong:} Daily operations are management's job. \\[2pt]
		\textbf{C wrong:} Workpapers are internal audit's. \\[2pt]
		\textbf{D wrong:} Hiring is management/HR.
	\end{frame}
	
	% =========================================================
	% Q45
	% =========================================================
	\begin{frame}[t]{Q45}
		\small
		\textbf{Question:} Which best describes "transparency" in governance?
		\begin{enumerate}[label=\Alph*.,itemsep=1pt,topsep=2pt]
			\item Timely and accurate disclosure of material information.
			\item Hiding information.
			\item Only disclosing good news.
			\item Disclosing only to management.
		\end{enumerate}
		\vspace{4pt}
		\textbf{Correct Answer: A}
	\end{frame}
	
	\begin{frame}[t]{Q45 -- Explanation}
		\small
		\textbf{Correct: A.} Transparency is timely and accurate disclosure of material information to stakeholders. \\[4pt]
		\textbf{B wrong:} Hiding is the opposite. \\[2pt]
		\textbf{C wrong:} Selective disclosure violates transparency. \\[2pt]
		\textbf{D wrong:} Disclosure is to stakeholders.
	\end{frame}
	
	% =========================================================
	% Q46
	% =========================================================
	\begin{frame}[t]{Q46}
		\small
		\textbf{Question:} Which best describes "accountability" in governance?
		\begin{enumerate}[label=\Alph*.,itemsep=1pt,topsep=2pt]
			\item Answerability for decisions and actions.
			\item Hiding decisions.
			\item Avoiding responsibility.
			\item Blaming others.
		\end{enumerate}
		\vspace{4pt}
		\textbf{Correct Answer: A}
	\end{frame}
	
	\begin{frame}[t]{Q46 -- Explanation}
		\small
		\textbf{Correct: A.} Accountability is answerability for decisions and actions. \\[4pt]
		\textbf{B wrong:} Hiding is the opposite. \\[2pt]
		\textbf{C wrong:} Avoiding responsibility is the opposite. \\[2pt]
		\textbf{D wrong:} Blaming others is not accountability.
	\end{frame}
	
	% =========================================================
	% Q47
	% =========================================================
	\begin{frame}[t]{Q47}
		\small
		\textbf{Question:} Which best describes "fairness" in governance?
		\begin{enumerate}[label=\Alph*.,itemsep=1pt,topsep=2pt]
			\item Equitable treatment of stakeholders.
			\item Favoring insiders.
			\item Ignoring minority shareholders.
			\item Hiding conflicts.
		\end{enumerate}
		\vspace{4pt}
		\textbf{Correct Answer: A}
	\end{frame}
	
	\begin{frame}[t]{Q47 -- Explanation}
		\small
		\textbf{Correct: A.} Fairness is equitable treatment of all stakeholders, including minority shareholders. \\[4pt]
		\textbf{B wrong:} Favoring insiders is the opposite. \\[2pt]
		\textbf{C wrong:} Ignoring minorities is the opposite. \\[2pt]
		\textbf{D wrong:} Hiding conflicts is unethical.
	\end{frame}
	
	% =========================================================
	% Q48
	% =========================================================
	\begin{frame}[t]{Q48}
		\small
		\textbf{Question:} Which best describes "responsibility" in governance?
		\begin{enumerate}[label=\Alph*.,itemsep=1pt,topsep=2pt]
			\item Acting ethically and in the best interest of the organization.
			\item Acting in self-interest.
			\item Ignoring stakeholders.
			\item Hiding information.
		\end{enumerate}
		\vspace{4pt}
		\textbf{Correct Answer: A}
	\end{frame}
	
	\begin{frame}[t]{Q48 -- Explanation}
		\small
		\textbf{Correct: A.} Responsibility is acting ethically and in the best interest of the organization and its stakeholders. \\[4pt]
		\textbf{B wrong:} Self-interest is not responsibility. \\[2pt]
		\textbf{C wrong:} Ignoring stakeholders is the opposite. \\[2pt]
		\textbf{D wrong:} Hiding is unethical.
	\end{frame}
	
	% =========================================================
	% Q49
	% =========================================================
	\begin{frame}[t]{Q49}
		\small
		\textbf{Question:} Which best describes "risk culture"?
		\begin{enumerate}[label=\Alph*.,itemsep=1pt,topsep=2pt]
			\item Shared values and behaviors regarding risk.
			\item The risk register.
			\item The audit plan.
			\item The marketing plan.
		\end{enumerate}
		\vspace{4pt}
		\textbf{Correct Answer: A}
	\end{frame}
	
	\begin{frame}[t]{Q49 -- Explanation}
		\small
		\textbf{Correct: A.} Risk culture is the shared values, beliefs, and behaviors regarding risk across the organization. \\[4pt]
		\textbf{B wrong:} The risk register is a tool. \\[2pt]
		\textbf{C wrong:} The audit plan is separate. \\[2pt]
		\textbf{D wrong:} Marketing is unrelated.
	\end{frame}
	
	% =========================================================
	% Q50
	% =========================================================
	\begin{frame}[t]{Q50}
		\small
		\textbf{Question:} Which best describes "risk appetite statement"?
		\begin{enumerate}[label=\Alph*.,itemsep=1pt,topsep=2pt]
			\item A formal statement of the amount/type of risk the entity is willing to accept.
			\item The audit plan.
			\item The marketing plan.
			\item The financial statements.
		\end{enumerate}
		\vspace{4pt}
		\textbf{Correct Answer: A}
	\end{frame}
	
	\begin{frame}[t]{Q50 -- Explanation}
		\small
		\textbf{Correct: A.} A risk appetite statement formally documents the amount and type of risk the entity is willing to accept. \\[4pt]
		\textbf{B wrong:} The audit plan is separate. \\[2pt]
		\textbf{C wrong:} Marketing is unrelated. \\[2pt]
		\textbf{D wrong:} Financial statements are outputs.
	\end{frame}
	
	% =========================================================
	% Q51
	% =========================================================
	\begin{frame}[t]{Q51}
		\small
		\textbf{Question:} Which is NOT a type of risk response?
		\begin{enumerate}[label=\Alph*.,itemsep=1pt,topsep=2pt]
			\item Avoid.
			\item Reduce.
			\item Ignore.
			\item Accept.
		\end{enumerate}
		\vspace{4pt}
		\textbf{Correct Answer: C}
	\end{frame}
	
	\begin{frame}[t]{Q51 -- Explanation}
		\small
		\textbf{Correct: C.} Ignoring is not a valid risk response. The four responses are avoid, reduce, share, accept. \\[4pt]
		\textbf{A wrong:} Avoid is a response. \\[2pt]
		\textbf{B wrong:} Reduce is a response. \\[2pt]
		\textbf{D wrong:} Accept is a response.
	\end{frame}
	
	% =========================================================
	% Q52
	% =========================================================
	\begin{frame}[t]{Q52}
		\small
		\textbf{Question:} Which best describes "risk owner"?
		\begin{enumerate}[label=\Alph*.,itemsep=1pt,topsep=2pt]
			\item The person accountable for managing a specific risk.
			\item The auditor.
			\item The board.
			\item The external regulator.
		\end{enumerate}
		\vspace{4pt}
		\textbf{Correct Answer: A}
	\end{frame}
	
	\begin{frame}[t]{Q52 -- Explanation}
		\small
		\textbf{Correct: A.} A risk owner is the person accountable for managing a specific risk. \\[4pt]
		\textbf{B wrong:} The auditor provides assurance. \\[2pt]
		\textbf{C wrong:} The board sets appetite. \\[2pt]
		\textbf{D wrong:} Regulators are external.
	\end{frame}
	
	% =========================================================
	% Q53
	% =========================================================
	\begin{frame}[t]{Q53}
		\small
		\textbf{Question:} Which best describes "risk map"?
		\begin{enumerate}[label=\Alph*.,itemsep=1pt,topsep=2pt]
			\item A visual representation of risks by likelihood and impact.
			\item The audit plan.
			\item The marketing plan.
			\item The financial statements.
		\end{enumerate}
		\vspace{4pt}
		\textbf{Correct Answer: A}
	\end{frame}
	
	\begin{frame}[t]{Q53 -- Explanation}
		\small
		\textbf{Correct: A.} A risk map plots risks by likelihood and impact to prioritize them. \\[4pt]
		\textbf{B wrong:} The audit plan is separate. \\[2pt]
		\textbf{C wrong:} Marketing is unrelated. \\[2pt]
		\textbf{D wrong:} Financial statements are outputs.
	\end{frame}
	
	% =========================================================
	% Q54
	% =========================================================
	\begin{frame}[t]{Q54}
		\small
		\textbf{Question:} Which best describes "key risk indicator" (KRI)?
		\begin{enumerate}[label=\Alph*.,itemsep=1pt,topsep=2pt]
			\item A metric that provides early warning of increasing risk.
			\item The audit plan.
			\item The marketing plan.
			\item The financial statements.
		\end{enumerate}
		\vspace{4pt}
		\textbf{Correct Answer: A}
	\end{frame}
	
	\begin{frame}[t]{Q54 -- Explanation}
		\small
		\textbf{Correct: A.} A KRI is a metric that provides early warning of increasing risk exposure. \\[4pt]
		\textbf{B wrong:} The audit plan is separate. \\[2pt]
		\textbf{C wrong:} Marketing is unrelated. \\[2pt]
		\textbf{D wrong:} Financial statements are outputs.
	\end{frame}
	
	% =========================================================
	% Q55
	% =========================================================
	\begin{frame}[t]{Q55}
		\small
		\textbf{Question:} Which best describes "risk-based internal auditing"?
		\begin{enumerate}[label=\Alph*.,itemsep=1pt,topsep=2pt]
			\item Aligning audit activities with the organization's risk profile.
			\item Auditing everything equally.
			\item Ignoring risk.
			\item Auditing only financial areas.
		\end{enumerate}
		\vspace{4pt}
		\textbf{Correct Answer: A}
	\end{frame}
	
	\begin{frame}[t]{Q55 -- Explanation}
		\small
		\textbf{Correct: A.} Risk-based auditing aligns audit activities with the organization's risk profile. \\[4pt]
		\textbf{B wrong:} Auditing everything equally is not risk-based. \\[2pt]
		\textbf{C wrong:} Ignoring risk is the opposite. \\[2pt]
		\textbf{D wrong:} Auditing only financial areas is too narrow.
	\end{frame}
	
	% =========================================================
	% Q56
	% =========================================================
	\begin{frame}[t]{Q56}
		\small
		\textbf{Question:} Which best describes "control self-assessment" (CSA)?
		\begin{enumerate}[label=\Alph*.,itemsep=1pt,topsep=2pt]
			\item A workshop-based approach where management assesses controls.
			\item An external audit technique.
			\item A type of fraud.
			\item A regulatory filing.
		\end{enumerate}
		\vspace{4pt}
		\textbf{Correct Answer: A}
	\end{frame}
	
	\begin{frame}[t]{Q56 -- Explanation}
		\small
		\textbf{Correct: A.} CSA engages management/process owners in assessing their own controls, often in facilitated workshops. \\[4pt]
		\textbf{B wrong:} CSA is internal. \\[2pt]
		\textbf{C wrong:} It is not fraud. \\[2pt]
		\textbf{D wrong:} It is not a filing.
	\end{frame}
	
	% =========================================================
	% Q57
	% =========================================================
	\begin{frame}[t]{Q57}
		\small
		\textbf{Question:} Which best describes "control gap"?
		\begin{enumerate}[label=\Alph*.,itemsep=1pt,topsep=2pt]
			\item A missing or inadequate control.
			\item A documented control.
			\item A management disagreement.
			\item Zero risk.
		\end{enumerate}
		\vspace{4pt}
		\textbf{Correct Answer: A}
	\end{frame}
	
	\begin{frame}[t]{Q57 -- Explanation}
		\small
		\textbf{Correct: A.} A control gap is a missing or inadequate control that leaves risk unmitigated. \\[4pt]
		\textbf{B wrong:} Documentation is a strength. \\[2pt]
		\textbf{C wrong:} Disagreement doesn't create a gap. \\[2pt]
		\textbf{D wrong:} Zero risk is unrealistic.
	\end{frame}
	
	% =========================================================
	% Q58
	% =========================================================
	\begin{frame}[t]{Q58}
		\small
		\textbf{Question:} Which best describes "compensating control"?
		\begin{enumerate}[label=\Alph*.,itemsep=1pt,topsep=2pt]
			\item A control that mitigates risk when the primary control is missing.
			\item A documented control.
			\item A management disagreement.
			\item Zero risk.
		\end{enumerate}
		\vspace{4pt}
		\textbf{Correct Answer: A}
	\end{frame}
	
	\begin{frame}[t]{Q58 -- Explanation}
		\small
		\textbf{Correct: A.} A compensating control mitigates risk when the primary control is missing or weak. \\[4pt]
		\textbf{B wrong:} Documentation is separate. \\[2pt]
		\textbf{C wrong:} Disagreement doesn't create a control. \\[2pt]
		\textbf{D wrong:} Zero risk is unrealistic.
	\end{frame}
	
	% =========================================================
	% Q59
	% =========================================================
	\begin{frame}[t]{Q59}
		\small
		\textbf{Question:} Which best describes "key control"?
		\begin{enumerate}[label=\Alph*.,itemsep=1pt,topsep=2pt]
			\item A control essential to mitigating a significant risk.
			\item A minor control.
			\item A management disagreement.
			\item Zero risk.
		\end{enumerate}
		\vspace{4pt}
		\textbf{Correct Answer: A}
	\end{frame}
	
	\begin{frame}[t]{Q59 -- Explanation}
		\small
		\textbf{Correct: A.} A key control is essential to mitigating a significant risk to an acceptable level. \\[4pt]
		\textbf{B wrong:} Minor controls are not key. \\[2pt]
		\textbf{C wrong:} Disagreement doesn't create a control. \\[2pt]
		\textbf{D wrong:} Zero risk is unrealistic.
	\end{frame}
	
	% =========================================================
	% Q60
	% =========================================================
	\begin{frame}[t]{Q60}
		\small
		\textbf{Question:} Which best describes "segregation of duties"?
		\begin{enumerate}[label=\Alph*.,itemsep=1pt,topsep=2pt]
			\item Separating authorization, custody, and record-keeping.
			\item Combining all duties in one person.
			\item Ignoring duties.
			\item Outsourcing.
		\end{enumerate}
		\vspace{4pt}
		\textbf{Correct Answer: A}
	\end{frame}
	
	\begin{frame}[t]{Q60 -- Explanation}
		\small
		\textbf{Correct: A.} Segregation of duties separates authorization, custody, and record-keeping to prevent fraud/error. \\[4pt]
		\textbf{B wrong:} Combining duties is the opposite. \\[2pt]
		\textbf{C wrong:} Ignoring is not a control. \\[2pt]
		\textbf{D wrong:} Outsourcing is a separate decision.
	\end{frame}
	
	% =========================================================
	% Q61
	% =========================================================
	\begin{frame}[t]{Q61}
		\small
		\textbf{Question:} Which best describes "IT general controls" (ITGC)?
		\begin{enumerate}[label=\Alph*.,itemsep=1pt,topsep=2pt]
			\item Controls over access, change management, and operations.
			\item Controls over financial reporting only.
			\item Controls over marketing.
			\item Controls over HR.
		\end{enumerate}
		\vspace{4pt}
		\textbf{Correct Answer: A}
	\end{frame}
	
	\begin{frame}[t]{Q61 -- Explanation}
		\small
		\textbf{Correct: A.} ITGCs cover access to programs/data, program changes, and computer operations. \\[4pt]
		\textbf{B wrong:} ITGCs are broader than financial reporting. \\[2pt]
		\textbf{C wrong:} Marketing is unrelated. \\[2pt]
		\textbf{D wrong:} HR is unrelated.
	\end{frame}
	
	% =========================================================
	% Q62
	% =========================================================
	\begin{frame}[t]{Q62}
		\small
		\textbf{Question:} Which best describes "application control"?
		\begin{enumerate}[label=\Alph*.,itemsep=1pt,topsep=2pt]
			\item A control embedded in an application to prevent/detect errors.
			\item A control over IT operations.
			\item A control over marketing.
			\item A control over HR.
		\end{enumerate}
		\vspace{4pt}
		\textbf{Correct Answer: A}
	\end{frame}
	
	\begin{frame}[t]{Q62 -- Explanation}
		\small
		\textbf{Correct: A.} Application controls are embedded in applications (e.g., input validation, edit checks) to prevent/detect errors. \\[4pt]
		\textbf{B wrong:} IT operations controls are ITGCs. \\[2pt]
		\textbf{C wrong:} Marketing is unrelated. \\[2pt]
		\textbf{D wrong:} HR is unrelated.
	\end{frame}
	
	% =========================================================
	% Q63
	% =========================================================
	\begin{frame}[t]{Q63}
		\small
		\textbf{Question:} Which best describes "ITGC access control"?
		\begin{enumerate}[label=\Alph*.,itemsep=1pt,topsep=2pt]
			\item Controls over who can access programs and data.
			\item Controls over marketing.
			\item Controls over HR.
			\item Controls over financial reporting.
		\end{enumerate}
		\vspace{4pt}
		\textbf{Correct Answer: A}
	\end{frame}
	
	\begin{frame}[t]{Q63 -- Explanation}
		\small
		\textbf{Correct: A.} Access controls manage who can access programs and data, and what they can do. \\[4pt]
		\textbf{B wrong:} Marketing is unrelated. \\[2pt]
		\textbf{C wrong:} HR is unrelated. \\[2pt]
		\textbf{D wrong:} Financial reporting is one area, not the definition.
	\end{frame}
	
	% =========================================================
	% Q64
	% =========================================================
	\begin{frame}[t]{Q64}
		\small
		\textbf{Question:} Which best describes "change management control"?
		\begin{enumerate}[label=\Alph*.,itemsep=1pt,topsep=2pt]
			\item Controls over changes to applications and systems.
			\item Controls over marketing.
			\item Controls over HR.
			\item Controls over financial reporting.
		\end{enumerate}
		\vspace{4pt}
		\textbf{Correct Answer: A}
	\end{frame}
	
	\begin{frame}[t]{Q64 -- Explanation}
		\small
		\textbf{Correct: A.} Change management controls ensure changes to applications/systems are authorized, tested, approved, and documented. \\[4pt]
		\textbf{B wrong:} Marketing is unrelated. \\[2pt]
		\textbf{C wrong:} HR is unrelated. \\[2pt]
		\textbf{D wrong:} Financial reporting is one area.
	\end{frame}
	
	% =========================================================
	% Q65
	% =========================================================
	\begin{frame}[t]{Q65}
		\small
		\textbf{Question:} Which best describes "IT operations control"?
		\begin{enumerate}[label=\Alph*.,itemsep=1pt,topsep=2pt]
			\item Controls over processing, backup, and recovery.
			\item Controls over marketing.
			\item Controls over HR.
			\item Controls over financial reporting.
		\end{enumerate}
		\vspace{4pt}
		\textbf{Correct Answer: A}
	\end{frame}
	
	\begin{frame}[t]{Q65 -- Explanation}
		\small
		\textbf{Correct: A.} IT operations controls cover processing schedules, backup, recovery, and incident management. \\[4pt]
		\textbf{B wrong:} Marketing is unrelated. \\[2pt]
		\textbf{C wrong:} HR is unrelated. \\[2pt]
		\textbf{D wrong:} Financial reporting is one area.
	\end{frame}
	
	% =========================================================
	% Q66
	% =========================================================
	\begin{frame}[t]{Q66}
		\small
		\textbf{Question:} Which best describes "governance failure"?
		\begin{enumerate}[label=\Alph*.,itemsep=1pt,topsep=2pt]
			\item Breakdown in oversight, accountability, or ethics.
			\item A minor control deficiency.
			\item A documentation gap.
			\item A training issue.
		\end{enumerate}
		\vspace{4pt}
		\textbf{Correct Answer: A}
	\end{frame}
	
	\begin{frame}[t]{Q66 -- Explanation}
		\small
		\textbf{Correct: A.} Governance failure is a breakdown in oversight, accountability, or ethics, often leading to scandals. \\[4pt]
		\textbf{B wrong:} Minor deficiencies are not governance failures. \\[2pt]
		\textbf{C wrong:} Documentation gaps are not governance failures. \\[2pt]
		\textbf{D wrong:} Training issues are not governance failures.
	\end{frame}
	
	% =========================================================
	% Q67
	% =========================================================
	\begin{frame}[t]{Q67}
		\small
		\textbf{Question:} Which best describes "whistleblowing" in governance?
		\begin{enumerate}[label=\Alph*.,itemsep=1pt,topsep=2pt]
			\item Reporting illegal or unethical activity through proper channels.
			\item Leaking to the media.
			\item Ignoring wrongdoing.
			\item Confronting the suspect.
		\end{enumerate}
		\vspace{4pt}
		\textbf{Correct Answer: A}
	\end{frame}
	
	\begin{frame}[t]{Q67 -- Explanation}
		\small
		\textbf{Correct: A.} Whistleblowing is reporting illegal/unethical activity through internal channels (or regulators if internal channels fail). \\[4pt]
		\textbf{B wrong:} Media leaks are a last resort. \\[2pt]
		\textbf{C wrong:} Ignoring is unethical. \\[2pt]
		\textbf{D wrong:} Confronting the suspect can compromise the investigation.
	\end{frame}
	
	% =========================================================
	% Q68
	% =========================================================
	\begin{frame}[t]{Q68}
		\small
		\textbf{Question:} Which best describes "ethics program"?
		\begin{enumerate}[label=\Alph*.,itemsep=1pt,topsep=2pt]
			\item A formal program to promote ethical behavior and compliance.
			\item A marketing campaign.
			\item A financial audit.
			\item A training session.
		\end{enumerate}
		\vspace{4pt}
		\textbf{Correct Answer: A}
	\end{frame}
	
	\begin{frame}[t]{Q68 -- Explanation}
		\small
		\textbf{Correct: A.} An ethics program promotes ethical behavior and compliance through codes, training, reporting, and enforcement. \\[4pt]
		\textbf{B wrong:} Marketing is unrelated. \\[2pt]
		\textbf{C wrong:} Financial audit is assurance. \\[2pt]
		\textbf{D wrong:} Training is one element.
	\end{frame}
	
	% =========================================================
	% Q69
	% =========================================================
	\begin{frame}[t]{Q69}
		\small
		\textbf{Question:} Which best describes "code of conduct"?
		\begin{enumerate}[label=\Alph*.,itemsep=1pt,topsep=2pt]
			\item A formal statement of expected behavior.
			\item A marketing plan.
			\item A financial audit.
			\item A training session.
		\end{enumerate}
		\vspace{4pt}
		\textbf{Correct Answer: A}
	\end{frame}
	
	\begin{frame}[t]{Q69 -- Explanation}
		\small
		\textbf{Correct: A.} A code of conduct formally states expected behavior for employees and stakeholders. \\[4pt]
		\textbf{B wrong:} Marketing is unrelated. \\[2pt]
		\textbf{C wrong:} Financial audit is assurance. \\[2pt]
		\textbf{D wrong:} Training is one element.
	\end{frame}
	
	% =========================================================
	% Q70
	% =========================================================
	\begin{frame}[t]{Q70}
		\small
		\textbf{Question:} Which best describes "tone at the top" in governance?
		\begin{enumerate}[label=\Alph*.,itemsep=1pt,topsep=2pt]
			\item Leadership's commitment to ethics and control.
			\item The CEO's salary.
			\item The marketing slogan.
			\item The IT password policy.
		\end{enumerate}
		\vspace{4pt}
		\textbf{Correct Answer: A}
	\end{frame}
	
	\begin{frame}[t]{Q70 -- Explanation}
		\small
		\textbf{Correct: A.} Tone at the top is leadership's commitment to ethics and control. \\[4pt]
		\textbf{B wrong:} Compensation is unrelated. \\[2pt]
		\textbf{C wrong:} Marketing slogans are not tone at the top. \\[2pt]
		\textbf{D wrong:} Password policy is one IT control.
	\end{frame}
	
	% =========================================================
	% Q71
	% =========================================================
	\begin{frame}[t]{Q71}
		\small
		\textbf{Question:} Which best describes "risk oversight"?
		\begin{enumerate}[label=\Alph*.,itemsep=1pt,topsep=2pt]
			\item The board's monitoring of management's risk management.
			\item The auditor's risk assessment.
			\item The marketing plan.
			\item The financial statements.
		\end{enumerate}
		\vspace{4pt}
		\textbf{Correct Answer: A}
	\end{frame}
	
	\begin{frame}[t]{Q71 -- Explanation}
		\small
		\textbf{Correct: A.} Risk oversight is the board's monitoring of management's risk management activities. \\[4pt]
		\textbf{B wrong:} The auditor's risk assessment is separate. \\[2pt]
		\textbf{C wrong:} Marketing is unrelated. \\[2pt]
		\textbf{D wrong:} Financial statements are outputs.
	\end{frame}
	
	% =========================================================
	% Q72
	% =========================================================
	\begin{frame}[t]{Q72}
		\small
		\textbf{Question:} Which best describes "risk governance"?
		\begin{enumerate}[label=\Alph*.,itemsep=1pt,topsep=2pt]
			\item The framework for how risk is managed and overseen.
			\item The audit plan.
			\item The marketing plan.
			\item The financial statements.
		\end{enumerate}
		\vspace{4pt}
		\textbf{Correct Answer: A}
	\end{frame}
	
	\begin{frame}[t]{Q72 -- Explanation}
		\small
		\textbf{Correct: A.} Risk governance is the framework for how risk is managed and overseen across the organization. \\[4pt]
		\textbf{B wrong:} The audit plan is separate. \\[2pt]
		\textbf{C wrong:} Marketing is unrelated. \\[2pt]
		\textbf{D wrong:} Financial statements are outputs.
	\end{frame}
	
	% =========================================================
	% Q73
	% =========================================================
	\begin{frame}[t]{Q73}
		\small
		\textbf{Question:} Which best describes "risk management framework"?
		\begin{enumerate}[label=\Alph*.,itemsep=1pt,topsep=2pt]
			\item A structured approach to identifying, assessing, and managing risk.
			\item The audit plan.
			\item The marketing plan.
			\item The financial statements.
		\end{enumerate}
		\vspace{4pt}
		\textbf{Correct Answer: A}
	\end{frame}
	
	\begin{frame}[t]{Q73 -- Explanation}
		\small
		\textbf{Correct: A.} A risk management framework is a structured approach to identifying, assessing, responding to, and monitoring risk. \\[4pt]
		\textbf{B wrong:} The audit plan is separate. \\[2pt]
		\textbf{C wrong:} Marketing is unrelated. \\[2pt]
		\textbf{D wrong:} Financial statements are outputs.
	\end{frame}
	
	% =========================================================
	% Q74
	% =========================================================
	\begin{frame}[t]{Q74}
		\small
		\textbf{Question:} Which best describes "COSO ERM 2017"?
		\begin{enumerate}[label=\Alph*.,itemsep=1pt,topsep=2pt]
			\item A framework integrating risk with strategy and performance.
			\item A financial reporting framework.
			\item A marketing framework.
			\item An HR framework.
		\end{enumerate}
		\vspace{4pt}
		\textbf{Correct Answer: A}
	\end{frame}
	
	\begin{frame}[t]{Q74 -- Explanation}
		\small
		\textbf{Correct: A.} COSO ERM 2017 integrates risk with strategy-setting and performance. \\[4pt]
		\textbf{B wrong:} Financial reporting is COSO IC. \\[2pt]
		\textbf{C wrong:} Marketing is unrelated. \\[2pt]
		\textbf{D wrong:} HR is unrelated.
	\end{frame}
	
	% =========================================================
	% Q75
	% =========================================================
	\begin{frame}[t]{Q75}
		\small
		\textbf{Question:} Which best describes "ISO 31000"?
		\begin{enumerate}[label=\Alph*.,itemsep=1pt,topsep=2pt]
			\item An international risk management standard.
			\item A financial reporting standard.
			\item A marketing standard.
			\item An HR standard.
		\end{enumerate}
		\vspace{4pt}
		\textbf{Correct Answer: A}
	\end{frame}
	
	\begin{frame}[t]{Q75 -- Explanation}
		\small
		\textbf{Correct: A.} ISO 31000 is an international standard providing principles and guidelines for risk management. \\[4pt]
		\textbf{B wrong:} Financial reporting is separate. \\[2pt]
		\textbf{C wrong:} Marketing is unrelated. \\[2pt]
		\textbf{D wrong:} HR is unrelated.
	\end{frame}
	
	% =========================================================
	% Q76
	% =========================================================
	\begin{frame}[t]{Q76}
		\small
		\textbf{Question:} Which best describes "risk-based audit plan"?
		\begin{enumerate}[label=\Alph*.,itemsep=1pt,topsep=2pt]
			\item An audit plan aligned with the organization's risk profile.
			\item An audit plan based on management requests.
			\item An audit plan based on prior findings.
			\item An audit plan based on budget.
		\end{enumerate}
		\vspace{4pt}
		\textbf{Correct Answer: A}
	\end{frame}
	
	\begin{frame}[t]{Q76 -- Explanation}
		\small
		\textbf{Correct: A.} A risk-based audit plan aligns audit activities with the organization's risk profile. \\[4pt]
		\textbf{B wrong:} Management requests are considered but not primary. \\[2pt]
		\textbf{C wrong:} Prior findings inform but risk drives. \\[2pt]
		\textbf{D wrong:} Budget is a constraint.
	\end{frame}
	
	% =========================================================
	% Q77
	% =========================================================
	\begin{frame}[t]{Q77}
		\small
		\textbf{Question:} Which best describes "assurance mapping"?
		\begin{enumerate}[label=\Alph*.,itemsep=1pt,topsep=2pt]
			\item Aligning assurance activities to risks and controls.
			\item Mapping the marketing plan.
			\item Mapping the HR plan.
			\item Mapping the financial statements.
		\end{enumerate}
		\vspace{4pt}
		\textbf{Correct Answer: A}
	\end{frame}
	
	\begin{frame}[t]{Q77 -- Explanation}
		\small
		\textbf{Correct: A.} Assurance mapping aligns assurance activities (internal audit, risk, compliance, external audit) to risks and controls. \\[4pt]
		\textbf{B wrong:} Marketing is unrelated. \\[2pt]
		\textbf{C wrong:} HR is unrelated. \\[2pt]
		\textbf{D wrong:} Financial statements are outputs.
	\end{frame}
	
	% =========================================================
	% Q78
	% =========================================================
	\begin{frame}[t]{Q78}
		\small
		\textbf{Question:} Which best describes "combined assurance"?
		\begin{enumerate}[label=\Alph*.,itemsep=1pt,topsep=2pt]
			\item Coordinating assurance from multiple providers.
			\item A single assurance provider.
			\item A marketing plan.
			\item A financial statement.
		\end{enumerate}
		\vspace{4pt}
		\textbf{Correct Answer: A}
	\end{frame}
	
	\begin{frame}[t]{Q78 -- Explanation}
		\small
		\textbf{Correct: A.} Combined assurance coordinates assurance from management, risk, compliance, internal audit, and external audit. \\[4pt]
		\textbf{B wrong:} Multiple providers are involved. \\[2pt]
		\textbf{C wrong:} Marketing is unrelated. \\[2pt]
		\textbf{D wrong:} Financial statements are outputs.
	\end{frame}
	
	% =========================================================
	% Q79
	% =========================================================
	\begin{frame}[t]{Q79}
		\small
		\textbf{Question:} Which best describes "governance, risk, and control" (GRC)?
		\begin{enumerate}[label=\Alph*.,itemsep=1pt,topsep=2pt]
			\item An integrated approach to governance, risk, and compliance.
			\item A marketing plan.
			\item A financial statement.
			\item An HR plan.
		\end{enumerate}
		\vspace{4pt}
		\textbf{Correct Answer: A}
	\end{frame}
	
	\begin{frame}[t]{Q79 -- Explanation}
		\small
		\textbf{Correct: A.} GRC is an integrated approach to governance, risk management, and compliance. \\[4pt]
		\textbf{B wrong:} Marketing is unrelated. \\[2pt]
		\textbf{C wrong:} Financial statements are outputs. \\[2pt]
		\textbf{D wrong:} HR is unrelated.
	\end{frame}
	
	% =========================================================
	% Q80
	% =========================================================
	\begin{frame}[t]{Q80}
		\small
		\textbf{Question:} Which best describes "risk-based internal audit"?
		\begin{enumerate}[label=\Alph*.,itemsep=1pt,topsep=2pt]
			\item Auditing areas of highest risk first.
			\item Auditing everything equally.
			\item Auditing only financial areas.
			\item Ignoring risk.
		\end{enumerate}
		\vspace{4pt}
		\textbf{Correct Answer: A}
	\end{frame}
	
	\begin{frame}[t]{Q80 -- Explanation}
		\small
		\textbf{Correct: A.} Risk-based internal audit focuses on areas of highest risk to the organization. \\[4pt]
		\textbf{B wrong:} Auditing everything equally is not risk-based. \\[2pt]
		\textbf{C wrong:} Auditing only financial is too narrow. \\[2pt]
		\textbf{D wrong:} Ignoring risk is the opposite.
	\end{frame}
	
	% =========================================================
	% Q81
	% =========================================================
	\begin{frame}[t]{Q81}
		\small
		\textbf{Question:} Which best describes "residual risk vs. inherent risk"?
		\begin{enumerate}[label=\Alph*.,itemsep=1pt,topsep=2pt]
			\item Residual = after controls; inherent = before controls.
			\item Residual = before controls; inherent = after controls.
			\item They are the same.
			\item Neither matters.
		\end{enumerate}
		\vspace{4pt}
		\textbf{Correct Answer: A}
	\end{frame}
	
	\begin{frame}[t]{Q81 -- Explanation}
		\small
		\textbf{Correct: A.} Residual risk is after controls; inherent risk is before controls. \\[4pt]
		\textbf{B wrong:} Definitions are reversed. \\[2pt]
		\textbf{C wrong:} They are distinct. \\[2pt]
		\textbf{D wrong:} Both matter.
	\end{frame}
	
	% =========================================================
	% Q82
	% =========================================================
	\begin{frame}[t]{Q82}
		\small
		\textbf{Question:} Which best describes "risk-based audit planning"?
		\begin{enumerate}[label=\Alph*.,itemsep=1pt,topsep=2pt]
			\item Prioritizing audit engagements based on risk.
			\item Prioritizing audits based on budget.
			\item Prioritizing audits based on management requests.
			\item Prioritizing audits based on prior findings.
		\end{enumerate}
		\vspace{4pt}
		\textbf{Correct Answer: A}
	\end{frame}
	
	\begin{frame}[t]{Q82 -- Explanation}
		\small
		\textbf{Correct: A.} Risk-based planning prioritizes engagements based on risk to the organization. \\[4pt]
		\textbf{B wrong:} Budget is a constraint. \\[2pt]
		\textbf{C wrong:} Management requests are considered. \\[2pt]
		\textbf{D wrong:} Prior findings inform but risk drives.
	\end{frame}
	
	% =========================================================
	% Q83
	% =========================================================
	\begin{frame}[t]{Q83}
		\small
		\textbf{Question:} Which best describes "governance framework"?
		\begin{enumerate}[label=\Alph*.,itemsep=1pt,topsep=2pt]
			\item A structure for directing and controlling the organization.
			\item The audit plan.
			\item The marketing plan.
			\item The financial statements.
		\end{enumerate}
		\vspace{4pt}
		\textbf{Correct Answer: A}
	\end{frame}
	
	\begin{frame}[t]{Q83 -- Explanation}
		\small
		\textbf{Correct: A.} A governance framework is the structure for directing, controlling, and holding the organization accountable. \\[4pt]
		\textbf{B wrong:} The audit plan is separate. \\[2pt]
		\textbf{C wrong:} Marketing is unrelated. \\[2pt]
		\textbf{D wrong:} Financial statements are outputs.
	\end{frame}
	
	% =========================================================
	% Q84
	% =========================================================
	\begin{frame}[t]{Q84}
		\small
		\textbf{Question:} Which best describes "board committee"?
		\begin{enumerate}[label=\Alph*.,itemsep=1pt,topsep=2pt]
			\item A sub-group of the board focused on specific oversight areas.
			\item A management committee.
			\item An audit team.
			\item A marketing team.
		\end{enumerate}
		\vspace{4pt}
		\textbf{Correct Answer: A}
	\end{frame}
	
	\begin{frame}[t]{Q84 -- Explanation}
		\small
		\textbf{Correct: A.} Board committees (audit, risk, compensation, nominating) focus on specific oversight areas. \\[4pt]
		\textbf{B wrong:} Management committees are separate. \\[2pt]
		\textbf{C wrong:} Audit teams are internal audit. \\[2pt]
		\textbf{D wrong:} Marketing teams are unrelated.
	\end{frame}
	
	% =========================================================
	% Q85
	% =========================================================
	\begin{frame}[t]{Q85}
		\small
		\textbf{Question:} Which best describes "audit committee charter"?
		\begin{enumerate}[label=\Alph*.,itemsep=1pt,topsep=2pt]
			\item A document defining the audit committee's purpose and responsibilities.
			\item The audit plan.
			\item The marketing plan.
			\item The financial statements.
		\end{enumerate}
		\vspace{4pt}
		\textbf{Correct Answer: A}
	\end{frame}
	
	\begin{frame}[t]{Q85 -- Explanation}
		\small
		\textbf{Correct: A.} The audit committee charter defines its purpose, authority, and responsibilities. \\[4pt]
		\textbf{B wrong:} The audit plan is separate. \\[2pt]
		\textbf{C wrong:} Marketing is unrelated. \\[2pt]
		\textbf{D wrong:} Financial statements are outputs.
	\end{frame}
	
	% =========================================================
	% Q86
	% =========================================================
	\begin{frame}[t]{Q86}
		\small
		\textbf{Question:} Which best describes "governance code"?
		\begin{enumerate}[label=\Alph*.,itemsep=1pt,topsep=2pt]
			\item A formal statement of governance principles and practices.
			\item The audit plan.
			\item The marketing plan.
			\item The financial statements.
		\end{enumerate}
		\vspace{4pt}
		\textbf{Correct Answer: A}
	\end{frame}
	
	\begin{frame}[t]{Q86 -- Explanation}
		\small
		\textbf{Correct: A.} A governance code formally states governance principles and practices. \\[4pt]
		\textbf{B wrong:} The audit plan is separate. \\[2pt]
		\textbf{C wrong:} Marketing is unrelated. \\[2pt]
		\textbf{D wrong:} Financial statements are outputs.
	\end{frame}
	
	% =========================================================
	% Q87
	% =========================================================
	\begin{frame}[t]{Q87}
		\small
		\textbf{Question:} Which best describes "stakeholder"?
		\begin{enumerate}[label=\Alph*.,itemsep=1pt,topsep=2pt]
			\item A party with an interest in the organization.
			\item Only shareholders.
			\item Only management.
			\item Only regulators.
		\end{enumerate}
		\vspace{4pt}
		\textbf{Correct Answer: A}
	\end{frame}
	
	\begin{frame}[t]{Q87 -- Explanation}
		\small
		\textbf{Correct: A.} Stakeholders include shareholders, employees, customers, suppliers, regulators, and communities. \\[4pt]
		\textbf{B wrong:} Shareholders are one group. \\[2pt]
		\textbf{C wrong:} Management is one group. \\[2pt]
		\textbf{D wrong:} Regulators are one group.
	\end{frame}
	
	% =========================================================
	% Q88
	% =========================================================
	\begin{frame}[t]{Q88}
		\small
		\textbf{Question:} Which best describes "shareholder vs. stakeholder"?
		\begin{enumerate}[label=\Alph*.,itemsep=1pt,topsep=2pt]
			\item Shareholders own equity; stakeholders have an interest.
			\item They are the same.
			\item Stakeholders own equity.
			\item Shareholders have no interest.
		\end{enumerate}
		\vspace{4pt}
		\textbf{Correct Answer: A}
	\end{frame}
	
	\begin{frame}[t]{Q88 -- Explanation}
		\small
		\textbf{Correct: A.} Shareholders own equity; stakeholders include anyone with an interest (employees, customers, etc.). \\[4pt]
		\textbf{B wrong:} They are distinct. \\[2pt]
		\textbf{C wrong:} Stakeholders don't necessarily own equity. \\[2pt]
		\textbf{D wrong:} Shareholders have an interest.
	\end{frame}
	
	% =========================================================
	% Q89
	% =========================================================
	\begin{frame}[t]{Q89}
		\small
		\textbf{Question:} Which best describes "agency problem"?
		\begin{enumerate}[label=\Alph*.,itemsep=1pt,topsep=2pt]
			\item Conflict of interest between principals (shareholders) and agents (management).
			\item A marketing problem.
			\item A financial statement error.
			\item An HR issue.
		\end{enumerate}
		\vspace{4pt}
		\textbf{Correct Answer: A}
	\end{frame}
	
	\begin{frame}[t]{Q89 -- Explanation}
		\small
		\textbf{Correct: A.} The agency problem is the conflict between principals (shareholders) and agents (management) due to differing interests. \\[4pt]
		\textbf{B wrong:} Marketing is unrelated. \\[2pt]
		\textbf{C wrong:} Financial statement errors are separate. \\[2pt]
		\textbf{D wrong:} HR is unrelated.
	\end{frame}
	
	% =========================================================
	% Q90
	% =========================================================
	\begin{frame}[t]{Q90}
		\small
		\textbf{Question:} Which best describes "monitoring" in the context of governance?
		\begin{enumerate}[label=\Alph*.,itemsep=1pt,topsep=2pt]
			\item Ongoing oversight of management and controls.
			\item Preparing the budget.
			\item Hiring staff.
			\item Setting the risk appetite.
		\end{enumerate}
		\vspace{4pt}
		\textbf{Correct Answer: A}
	\end{frame}
	
	\begin{frame}[t]{Q90 -- Explanation}
		\small
		\textbf{Correct: A.} Monitoring is ongoing oversight of management and controls by the board and others. \\[4pt]
		\textbf{B wrong:} Budgeting is a control activity. \\[2pt]
		\textbf{C wrong:} Hiring is a control activity. \\[2pt]
		\textbf{D wrong:} Risk appetite is set by the board.
	\end{frame}
	
	% =========================================================
	% Q91
	% =========================================================
	\begin{frame}[t]{Q91}
		\small
		\textbf{Question:} Which best describes "governance and risk" linkage?
		\begin{enumerate}[label=\Alph*.,itemsep=1pt,topsep=2pt]
			\item Governance sets the tone and oversight for risk management.
			\item Governance is separate from risk.
			\item Risk is only the auditor's job.
			\item Governance is only the board's job.
		\end{enumerate}
		\vspace{4pt}
		\textbf{Correct Answer: A}
	\end{frame}
	
	\begin{frame}[t]{Q91 -- Explanation}
		\small
		\textbf{Correct: A.} Governance sets the tone and oversight for risk management, including appetite and accountability. \\[4pt]
		\textbf{B wrong:} They are linked. \\[2pt]
		\textbf{C wrong:} Risk is management's job. \\[2pt]
		\textbf{D wrong:} Governance involves the board and management.
	\end{frame}
	
	% =========================================================
	% Q92
	% =========================================================
	\begin{frame}[t]{Q92}
		\small
		\textbf{Question:} Which best describes "control environment" vs. "control activities"?
		\begin{enumerate}[label=\Alph*.,itemsep=1pt,topsep=2pt]
			\item Environment = tone; activities = policies/procedures.
			\item Environment = policies; activities = tone.
			\item They are the same.
			\item Neither matters.
		\end{enumerate}
		\vspace{4pt}
		\textbf{Correct Answer: A}
	\end{frame}
	
	\begin{frame}[t]{Q92 -- Explanation}
		\small
		\textbf{Correct: A.} Control environment is the tone at the top; control activities are policies/procedures. \\[4pt]
		\textbf{B wrong:} Definitions are reversed. \\[2pt]
		\textbf{C wrong:} They are distinct. \\[2pt]
		\textbf{D wrong:} Both matter.
	\end{frame}
	
	% =========================================================
	% Q93
	% =========================================================
	\begin{frame}[t]{Q93}
		\small
		\textbf{Question:} Which best describes "risk assessment" vs. "risk response"?
		\begin{enumerate}[label=\Alph*.,itemsep=1pt,topsep=2pt]
			\item Assessment = identify/analyze; response = avoid/reduce/share/accept.
			\item Assessment = response; response = assessment.
			\item They are the same.
			\item Neither matters.
		\end{enumerate}
		\vspace{4pt}
		\textbf{Correct Answer: A}
	\end{frame}
	
	\begin{frame}[t]{Q93 -- Explanation}
		\small
		\textbf{Correct: A.} Assessment identifies and analyzes risk; response selects a treatment (avoid, reduce, share, accept). \\[4pt]
		\textbf{B wrong:} Definitions are reversed. \\[2pt]
		\textbf{C wrong:} They are distinct. \\[2pt]
		\textbf{D wrong:} Both matter.
	\end{frame}
	
	% =========================================================
	% Q94
	% =========================================================
	\begin{frame}[t]{Q94}
		\small
		\textbf{Question:} Which best describes "control" vs. "risk"?
		\begin{enumerate}[label=\Alph*.,itemsep=1pt,topsep=2pt]
			\item Control mitigates risk.
			\item Risk mitigates control.
			\item They are the same.
			\item Neither matters.
		\end{enumerate}
		\vspace{4pt}
		\textbf{Correct Answer: A}
	\end{frame}
	
	\begin{frame}[t]{Q94 -- Explanation}
		\small
		\textbf{Correct: A.} Controls mitigate risk to an acceptable level. \\[4pt]
		\textbf{B wrong:} Risk does not mitigate control. \\[2pt]
		\textbf{C wrong:} They are distinct. \\[2pt]
		\textbf{D wrong:} Both matter.
	\end{frame}
	
	% =========================================================
	% Q95
	% =========================================================
	\begin{frame}[t]{Q95}
		\small
		\textbf{Question:} Which best describes "governance" vs. "management"?
		\begin{enumerate}[label=\Alph*.,itemsep=1pt,topsep=2pt]
			\item Governance = oversight; management = execution.
			\item Governance = execution; management = oversight.
			\item They are the same.
			\item Neither matters.
		\end{enumerate}
		\vspace{4pt}
		\textbf{Correct Answer: A}
	\end{frame}
	
	\begin{frame}[t]{Q95 -- Explanation}
		\small
		\textbf{Correct: A.} Governance provides oversight; management executes operations. \\[4pt]
		\textbf{B wrong:} Roles are reversed. \\[2pt]
		\textbf{C wrong:} They are distinct. \\[2pt]
		\textbf{D wrong:} Both matter.
	\end{frame}
	
	% =========================================================
	% Q96
	% =========================================================
	\begin{frame}[t]{Q96}
		\small
		\textbf{Question:} Which best describes "risk appetite" vs. "risk tolerance"?
		\begin{enumerate}[label=\Alph*.,itemsep=1pt,topsep=2pt]
			\item Appetite = willing to accept; tolerance = acceptable variation.
			\item Appetite = variation; tolerance = willingness.
			\item They are the same.
			\item Neither matters.
		\end{enumerate}
		\vspace{4pt}
		\textbf{Correct Answer: A}
	\end{frame}
	
	\begin{frame}[t]{Q96 -- Explanation}
		\small
		\textbf{Correct: A.} Risk appetite is the amount/type of risk willing to accept; risk tolerance is the acceptable variation around it. \\[4pt]
		\textbf{B wrong:} Definitions are reversed. \\[2pt]
		\textbf{C wrong:} They are distinct. \\[2pt]
		\textbf{D wrong:} Both matter.
	\end{frame}
	
	% =========================================================
	% Q97
	% =========================================================
	\begin{frame}[t]{Q97}
		\small
		\textbf{Question:} Which best describes "inherent risk" vs. "residual risk"?
		\begin{enumerate}[label=\Alph*.,itemsep=1pt,topsep=2pt]
			\item Inherent = before controls; residual = after controls.
			\item Inherent = after controls; residual = before controls.
			\item They are the same.
			\item Neither matters.
		\end{enumerate}
		\vspace{4pt}
		\textbf{Correct Answer: A}
	\end{frame}
	
	\begin{frame}[t]{Q97 -- Explanation}
		\small
		\textbf{Correct: A.} Inherent risk is before controls; residual risk is after controls. \\[4pt]
		\textbf{B wrong:} Definitions are reversed. \\[2pt]
		\textbf{C wrong:} They are distinct. \\[2pt]
		\textbf{D wrong:} Both matter.
	\end{frame}
	
	% =========================================================
	% Q98
	% =========================================================
	\begin{frame}[t]{Q98}
		\small
		\textbf{Question:} Which best describes "control risk" vs. "detection risk"?
		\begin{enumerate}[label=\Alph*.,itemsep=1pt,topsep=2pt]
			\item Control risk = control fails; detection risk = auditor fails to detect.
			\item Control risk = auditor fails; detection risk = control fails.
			\item They are the same.
			\item Neither matters.
		\end{enumerate}
		\vspace{4pt}
		\textbf{Correct Answer: A}
	\end{frame}
	
	\begin{frame}[t]{Q98 -- Explanation}
		\small
		\textbf{Correct: A.} Control risk is the risk a control fails; detection risk is the risk the auditor fails to detect a misstatement. \\[4pt]
		\textbf{B wrong:} Definitions are reversed. \\[2pt]
		\textbf{C wrong:} They are distinct. \\[2pt]
		\textbf{D wrong:} Both matter.
	\end{frame}
	
	% =========================================================
	% Q99
	% =========================================================
	\begin{frame}[t]{Q99}
		\small
		\textbf{Question:} Which best describes "preventive" vs. "detective" controls?
		\begin{enumerate}[label=\Alph*.,itemsep=1pt,topsep=2pt]
			\item Preventive = stop; detective = find.
			\item Preventive = find; detective = stop.
			\item They are the same.
			\item Neither matters.
		\end{enumerate}
		\vspace{4pt}
		\textbf{Correct Answer: A}
	\end{frame}
	
	\begin{frame}[t]{Q99 -- Explanation}
		\small
		\textbf{Correct: A.} Preventive controls stop errors/fraud; detective controls find them after they occur. \\[4pt]
		\textbf{B wrong:} Definitions are reversed. \\[2pt]
		\textbf{C wrong:} They are distinct. \\[2pt]
		\textbf{D wrong:} Both matter.
	\end{frame}
	
	% =========================================================
	% Q100
	% =========================================================
	\begin{frame}[t]{Q100}
		\small
		\textbf{Question:} Which of the following best summarizes governance, risk management, and control?
		\begin{enumerate}[label=\Alph*.,itemsep=1pt,topsep=2pt]
			\item Governance sets direction and oversight; risk management identifies and treats risk; control mitigates risk.
			\item Governance executes; risk management ignores; control does nothing.
			\item Governance is optional; risk management is optional; control is optional.
			\item Governance is only for public companies.
		\end{enumerate}
		\vspace{4pt}
		\textbf{Correct Answer: A}
	\end{frame}
	
	\begin{frame}[t]{Q100 -- Explanation}
		\small
		\textbf{Correct: A.} Governance sets direction and oversight; risk management identifies and treats risk; controls mitigate risk to acceptable levels. \\[4pt]
		\textbf{B wrong:} The roles are misrepresented. \\[2pt]
		\textbf{C wrong:} All three are essential. \\[2pt]
		\textbf{D wrong:} All organizations need governance.
	\end{frame}
	
	\begin{frame}{End}
		\centering
		\Large End of 100-Question Governance, Risk Management, and Control Deck\\[8pt]
		\normalsize Review weak areas and re-test.\\[4pt]
		\small Add more questions by copying the Q/Explanation frame pairs.
	\end{frame}
	
\end{document}